\documentclass[a4paper]{jsarticle}

\usepackage[top=30truemm,bottom=27truemm,left=18truemm,right=18truemm,columnsep=7truemm]{geometry}
\usepackage[dvipdfmx]{graphicx,xcolor}
\usepackage[haranoaji]{pxchfon}
\usepackage{newtxtext,newtxmath}
\usepackage{booktabs}
\usepackage{array}
\usepackage{enumitem}
\usepackage{titlesec}
\usepackage{tikz}
\usepackage{url}
\usetikzlibrary{arrows.meta,positioning,shapes.geometric,fit}

\setlength{\columnsep}{7truemm}
\setlength{\parindent}{1zw}
\setlength{\parskip}{0pt}
\setlength{\textfloatsep}{7pt plus 2pt minus 2pt}
\setlength{\floatsep}{6pt plus 2pt minus 2pt}
\setlength{\intextsep}{6pt plus 2pt minus 2pt}
\setlist[itemize]{leftmargin=1.5zw,itemsep=0pt,topsep=2pt,parsep=0pt}
\pagestyle{empty}

\makeatletter
\renewcommand{\footnoterule}{}
\renewcommand{\@makefntext}[1]{\noindent #1}
\makeatother

\titleformat{\section}
  {\normalfont\fontsize{10.5pt}{16pt}\selectfont\bfseries}
  {\thesection.}{0.6zw}{}
\titlespacing*{\section}{0pt}{8pt plus 2pt minus 1pt}{2pt}
\titleformat{\subsection}
  {\normalfont\fontsize{10.5pt}{16pt}\selectfont\bfseries}
  {\thesubsection}{0.6zw}{}
\titlespacing*{\subsection}{0pt}{6pt plus 1pt minus 1pt}{1pt}

\newcommand{\bodyfont}{\fontsize{10.5pt}{16pt}\selectfont}
\newcommand{\titlefont}{\fontsize{12pt}{18pt}\selectfont}
\newcommand{\authorfont}{\fontsize{10.5pt}{16pt}\selectfont}
\newcommand{\term}[1]{\textbf{#1}}

\newcolumntype{L}[1]{>{\raggedright\arraybackslash}p{#1}}
\newcolumntype{C}[1]{>{\centering\arraybackslash}p{#1}}

\makeatletter
\long\def\@makecaption#1#2{{\bodyfont
  \advance\leftskip .0628\linewidth
  \advance\rightskip .0628\linewidth
  \vskip\abovecaptionskip
  \sbox\@tempboxa{#1\hskip1zw\relax #2}%
  \ifdim \wd\@tempboxa <\hsize \centering \fi
  #1{\hskip1zw\relax}#2\par
  \vskip\belowcaptionskip}}
\makeatother

\begin{document}
\bodyfont

\twocolumn[
\begin{@twocolumnfalse}
  \begin{center}
    {\titlefont\bfseries 電子カルテ利用時の認知負荷軽減に向けたタスクグラフによるUI生成\par}
    \vspace{2pt}
    {\titlefont Task-Graph-Based EHR Interface Generation Toward Reducing Cognitive Load\par}
    \vspace{7pt}
    {\authorfont 平田 蓮\par}
    {\authorfont Lenne Hirata\par}
  \end{center}
  \vspace{8pt}
  \begin{center}
  \begin{minipage}{160truemm}
  \noindent\textbf{Abstract}\quad
  Electronic health records contain data needed for care, but clinicians often have to find and combine information across multiple screens. This research aims to reduce the cognitive load of that work by generating task-specific interfaces from a task graph. The graph connects a clinical scenario, information-checking tasks, data queries, and widgets. The system executes read-only SQL inside the hospital environment and renders an interactive interface without sending clinical data to the language model. It has been deployed on a hospital server, where its application and health endpoints were verified. A completed technical evaluation showed that cross-layer contracts rejected all 763 injected inconsistencies while accepting all 763 valid updates. This result establishes a prerequisite for evaluating use; it does not demonstrate reduced cognitive load. Subject to ethics review, institutional permission, and expert validation, we plan to compare the proposed interface with the current EHR before graduation and measure task success, completion time, interactions, NASA-TLX, SUS, and interview responses.\par
  \end{minipage}
  \end{center}
  \vspace{11pt}
\end{@twocolumnfalse}
]

\begingroup
\renewcommand{\thefootnote}{}
\footnotetext{\authorfont
  医療情報学講座・医療情報学分野 (指導教員: 黒田 知宏)}
\endgroup

\section{はじめに}

電子カルテには、患者属性、検査、処方、処置、診療記録など、診療に必要な情報が蓄積されている。一方、これらの情報は複数の表や画面に分かれている。医療従事者は診療の目的に応じて画面を移動し、必要な項目を探し、時間的な変化や情報同士の関係を頭の中で組み立てる。この情報探索と統合には、認知的な負担が伴う。

本研究の目的は、診療タスクに必要な情報を電子カルテから取り出し、確認しやすい画面にまとめ、情報探索と統合に伴う認知的な負担を軽減することである。同じ患者でも、術前診察ではアレルギー、既往歴、検査結果、処方を広く確認し、慢性疾患の外来では血圧や検査値の推移を追う。固定された画面を増やすだけでは、この違いに柔軟に対応しにくい。そこで、診療場面、行うタスク、参照するデータ、表示方法の関係を\term{タスクグラフ}として表し、グラフからタスク別のUIを生成する。

研究は三つの段階からなる。第一に、タスクグラフから電子カルテUIを生成するシステムを実装し、院内環境で動作させる。第二に、タスク、SQL、実行結果、表示の対応が壊れていないかを技術評価で確かめる。第三に、現行電子カルテと提案UIを同じ情報確認課題で比較し、作業成績と主観的な負荷を測る。現在までに、院内サーバーへの配置と起動確認、および整合性の技術評価を終えた。必要な専門家確認、倫理審査、院内の実施許可が整えば、卒業までにユーザー評価を行う。

本稿では、研究全体の目的とシステムを説明した後、完了した技術評価をその一部として報告する。さらに、認知負荷の軽減を検証するユーザー評価と、卒業までに作成する成果物を示す。

\section{関連研究}

\subsection{電子カルテの表示と認知負荷}

電子カルテの使いやすさは、利用者が目的を達成できる正確さ、効率、満足度とともに評価する必要がある。NISTの電子カルテ評価手順も、臨床専門家による確認と、代表的な利用者が現実的な課題を行う評価を組み合わせている\cite{lowry2012ehr}。したがって、提案UIの価値は画面の印象だけでなく、同じ診療課題をどの程度正確かつ少ない負担で完了できるかによって判断する。

先行研究では、情報を目的に合わせて整理した表示が、電子カルテでの情報探索を支援することが報告されている。入院患者の安全確認用ダッシュボードは、従来の電子カルテと比べて作業時間やマウス操作を減らした\cite{fuller2020dashboard}。問題ごとに情報をまとめた表示でも、標準画面より短い時間と少ない誤りで情報を取得でき、主観的な作業負荷も低下した\cite{semanik2021problem}。また、臨床的な概念に沿う可視化は、患者の優先順位付けに伴う医師の認知負荷と関係する\cite{pollack2020visualization}。これらはタスクに合う表示の可能性を示すが、利用者のタスクから表示を生成する方法は扱っていない。

\subsection{タスクに基づくUI生成}

生成UIでは、利用者の要求から画面を作り、生成後も要求と画面の関係を保つ必要がある。Jellyは利用者の情報タスクをデータモデルとして保持し、自然言語による指示と直接操作をモデルの変更へ変換する\cite{cao2025jelly}。

本研究は、診療場面から表示までをタスクグラフで結び、電子カルテのデータ取得まで含めてUIを生成する。臨床データを扱うため、UIを生成できることに加えて、要求した患者、項目、期間、情報源が表示まで保たれているかを検査する。

\section{研究全体の問いと評価}

最終的に確かめたいことは、タスクに合わせて生成したUIが、現行電子カルテで情報を探して統合する負担を減らせるかである。この問いに段階的に答えるため、次の三点を調べる。

\begin{itemize}
  \item RQ1：診療タスクと必要情報をタスクグラフで表し、院内の電子カルテデータから操作可能なUIを生成できるか。
  \item RQ2：タスクグラフ、SQL、実行結果、表示の対応を検査すると、想定した8種類の単一不整合をどこまで検出し、修復できるか。
  \item RQ3：提案UIは現行電子カルテと比べて、情報確認の正確さを損なわず、効率を高めて利用者の認知負荷を軽減できるか。
\end{itemize}

RQ1にはシステムの実装と院内動作、RQ2には完了した層間整合性評価を対応させる。RQ3は、必要な確認と許可が整った後に行うユーザー評価で検証する。表\ref{tab:research-map}に研究全体の流れを示す。

\begin{table*}[t]
  \centering
  \setlength{\tabcolsep}{4pt}
  \caption{研究全体における評価の位置付け}
  \label{tab:research-map}
  \begin{tabular}{C{0.07\textwidth}L{0.23\textwidth}L{0.28\textwidth}L{0.31\textwidth}}
    \toprule
    問い & 確認する内容 & 方法 & 判断に使う結果 \\
    \midrule
    RQ1 & タスクからUIを生成できるか & タスクグラフの実装と院内環境での動作確認 & グラフ検証、読み取り専用SQL、UI描画、院内サーバーでの起動 \\
    RQ2 & 要求と表示の対応を保てるか & 8種類の不整合を入れる技術評価 & 不整合の流出、妥当な更新の受理、問題箇所の特定、修復 \\
    RQ3 & 情報探索と統合の負担を減らせるか & 現行電子カルテと提案UIのユーザー評価 & タスク成功、完了時間、操作数、NASA Task Load Index、System Usability Scale、聞き取り \\
    \bottomrule
  \end{tabular}
\end{table*}

\section{タスクグラフによるUI生成}

\subsection{タスク、データ、表示の関係}

タスクグラフは、Scenario、Task、Data、Widgetの四つの要素からなる。Scenarioは術前診察や慢性疾患外来などの診療場面を表す。Taskには、その場面で利用者が確認する内容を記録する。Dataには参照する院内データと読み取り専用SQLを、Widgetには表、数値、時系列などの表示方法を記録する。TaskからDataとWidgetをたどることで、何のために、どのデータを、どの形で示すかを確認できる。

たとえば、術前診察で腎機能を確認するTaskには、クレアチニンなどを取得するDataと、直近値を示す表または推移を示すグラフを結び付ける。慢性疾患外来で血圧の変化を確認する場合は、同じ画面構造を固定せず、期間と表示方法をタスクに合わせて変える。データの所在を起点にせず、利用者が行う確認からUIを組み立てる点が特徴である。

\subsection{生成と実行の流れ}

図\ref{fig:ui-generation}に処理の流れを示す。利用者が文章で診療タスクを入力するか、既存の画面を操作して変更を指示すると、生成モデルが構造化されたタスクグラフを作る。システムはグラフの形式と参照関係を検証し、DataにあるSELECT文だけを院内のローカルデータベースに対して実行する。取得結果はWidgetへ渡し、タスクごとに並べた対話的な画面として表示する。

生成モデルへ送るのはタスクとデータ構造の情報であり、電子カルテから取得した患者データは送らない。患者データの取得と画面描画は院内環境で完結する。利用者が項目、期間、表示方法を変更したときも、変更要求をグラフへ反映してからデータを再取得する。

\begin{figure*}[t]
  \centering
  \resizebox{0.94\textwidth}{!}{%
  \begin{tikzpicture}[
    node distance=5mm and 7mm,
    box/.style={draw,rounded corners=1.2mm,minimum width=25mm,minimum height=11mm,align=center,fill=gray!5},
    graph/.style={draw,rounded corners=1.2mm,minimum width=18mm,minimum height=9mm,align=center,fill=blue!6},
    arr/.style={-{Latex[length=2.3mm]},thick}
  ]
    \node[box] (input) {利用者のタスク\\文章・画面操作};
    \node[graph,right=of input] (scenario) {Scenario\\診療場面};
    \node[graph,right=of scenario] (task) {Task\\確認内容};
    \node[graph,right=of task] (data) {Data\\参照先・SQL};
    \node[graph,right=of data] (widget) {Widget\\表示方法};
    \node[box,right=of widget] (ui) {タスク別の\\対話的UI};
    \node[box,below=8mm of data] (ehr) {院内データベース\\読み取り専用};
    \draw[arr] (input) -- node[above]{構造化} (scenario);
    \draw[arr] (scenario) -- (task);
    \draw[arr] (task) -- (data);
    \draw[arr] (data) -- (widget);
    \draw[arr] (widget) -- node[above]{描画} (ui);
    \draw[arr] (data) -- node[right]{SELECT} (ehr);
    \draw[arr] (ehr) -| node[below right]{取得結果} (widget);
    \draw[arr,dashed] (ui.north) to[bend right=22] node[above]{変更} (input.north);
    \node[draw,dashed,rounded corners,fit=(scenario)(task)(data)(widget),inner sep=3mm,label=below:{タスクグラフ}] {};
  \end{tikzpicture}}
  \caption{タスクグラフから電子カルテUIを生成する流れ}
  \label{fig:ui-generation}
\end{figure*}

\section{院内環境への実装}

提案システムをLinux向けのDockerイメージとして院内サーバーへ配置した。院内のデータウェアハウス（DWH）から取得した135表をローカルのSQLiteデータベースへ取り込み、モデルが要求する表名と列名の存在、データベースの整合性、ホストとコンテナ間のファイル一致を確認した。2026年8月24日の動作確認では、院内サーバー上で画面とヘルス確認の応答がともにHTTP 200となった。アプリと患者データは院内環境に置き、患者単位の情報をGitHubやNotionへ保存しない。

この動作確認により、RQ1のうち、院内データを使った読み取り専用SQLの実行と、取得結果のUI描画までを確認した。タスクの内容と表示の医学的な妥当性は、まだ確認していない。

院内実装により、運用中の電子カルテと同じ院内環境で比較するためのシステム面の準備が整った。ユーザー評価では、専門家が確認した同じ情報確認課題を、現行電子カルテと提案UIの双方で行う。提案UIが速く見える課題だけを選ばないように、必要情報、開始状態、完了条件、利用できる機能を事前にそろえる。

現時点で確認できたのは、院内データを使って提案システムが動作することまでである。現行電子カルテより認知負荷が低いか、情報を正確に取得できるかは、後述するユーザー評価で確かめる。

\section{整合性の技術評価}

\subsection{研究全体での役割}

ユーザー評価に使う画面が、要求と異なる患者、項目、期間を表示していれば、操作時間や認知負荷を測っても提案UIの効果を判断できない。そこで、利用評価に先立ち、タスクから表示までの対応を検査する技術評価を行った。この評価は研究全体の一部であり、認知負荷の軽減を直接測るものではない。

臨床質問と正解SQLにはEHRSQL-2024\cite{lee2024ehrsql}を用い、MIMIC-IV Clinical Database Demo v2.2\cite{johnson2023mimic}でSQLを実行した。基準となるタスク、SQL、実行結果、表示に、対象患者、対象項目、期間、集約方法、情報源、表示方法、データと表示の接続、更新後に残る古い結果という8種類の不整合を一つずつ入れた。

同じ候補を、個々のSQLや表示だけを調べる局所検査、成果物ごとの条件を調べる検査、タスクから表示までの関係を照合するグラフ契約、同じ関係を一枚の表に保持するサイドカー契約へ入力した。後者の二方式に同じ確認内容を持たせ、グラフという保存形式と、関係をまとめて検査することの効果を分けた。

\subsection{正式評価の結果}

評価条件を固定した後、最終評価用データのうち正解SQLを持つ934件を一度だけ実行した。全件で基準UIを構築でき、134種類の質問テンプレートから763組の更新前後ペアを作った。8種類の不整合と妥当な更新を各組に適用し、合計6,104回検証した。主な結果を表\ref{tab:integrity-results}に示す。

\begin{table}[t]
  \centering
  \setlength{\tabcolsep}{2pt}
  \caption{層間整合性の正式技術評価}
  \label{tab:integrity-results}
  \begin{tabular}{L{0.34\columnwidth}C{0.19\columnwidth}C{0.19\columnwidth}C{0.15\columnwidth}}
    \toprule
    検査方法 & 不整合の流出 & 妥当な更新の受理 & 特定・修復 \\
    \midrule
    局所検査 & 100.00\% & 100.00\% & --- \\
    成果物ごとの契約 & 48.10\% & 100.00\% & 一部 \\
    グラフ契約 & 0.00\% & 100.00\% & 763/763 \\
    サイドカー契約 & 0.00\% & 100.00\% & 763/763 \\
    \bottomrule
  \end{tabular}
\end{table}

局所検査は注入した不整合763件をすべて受理し、成果物ごとの契約も367件を受理した。グラフ契約とサイドカー契約は不整合をすべて拒否し、妥当な更新はすべて受理した。両方式は問題箇所の特定と一回の修復にも全件で成功し、判定差はなかった。したがって、不整合を防いだ要因はグラフ形式そのものではなく、タスク、データ、表示の関係を一つの契約として検査したことだと解釈する。

この結果により、ユーザー評価へ進む前提となる技術的な確認方法を得た。ただし、本評価は公開デモデータを使った決定的な評価であり、表示内容の医学的な妥当性、使いやすさ、認知負荷、臨床判断への影響は示していない。

\section{卒業までに行うユーザー評価}

\subsection{比較する条件}

ユーザー評価の最初の対象候補は、院内実装と既存のタスクモデルがある麻酔科術前外来である。評価を始める前に、診療専門家とともに、確認課題、正解、採点基準、必要情報、重大な見落とし、課題の終了条件を定める。提案UIと現行電子カルテでは、同じ目的と必要情報を持つ課題を用いる。参加者内で両方のシステムを使い、症例と提示順を入れ替えて順序の影響を抑える。

提案UIでは、専門家が確認したタスクグラフと画面を使用する。参加者へ提示する前に層間整合性検査を適用し、要求、SQL、実行結果、表示の対応を確認する。現行電子カルテでは、通常の操作経路から同じ情報を確認する。両条件で開始状態、利用できる情報、説明時間をそろえ、操作ログにはクリック、画面遷移、スクロール、完了時刻を記録する。

\subsection{評価指標}

認知負荷の軽減だけで判断せず、正確さ、効率、主観評価を分けて測る。主な指標は次のとおりである。

\begin{itemize}
  \item タスク成功率と、見落としや誤答の内容
  \item 成功した試行の完了時間、クリック、画面遷移、スクロール
  \item 作業負荷を測るNASA-TLX
  \item 使いやすさを測るSystem Usability Scale（SUS）
  \item 情報の探し方、迷った箇所、表示への信頼に関する聞き取り
\end{itemize}

タスク成功率を保ったうえで、完了時間や操作が減り、NASA-TLXが低下するかを中心に検討する。SUSと聞き取りは、結果の理由と改善点を解釈するために用いる。対象者の条件と人数、割付、除外基準、欠測の扱い、主要指標、分析方法は、予備確認後かつ本評価の前に固定する。

\subsection{倫理と安全}

人を対象とする評価は、倫理審査と院内の実施許可、募集、同意、撤回、ログと質問紙の保存方法を確認してから始める。診療中の判断には用いず、専門家が確認した合成症例と課題を使う。実患者情報を研究用成果物、Notion、GitHubへ保存しない。専門家確認と許可が整わない場合は参加者を募集せず、準備状況と未実施の理由を修士論文へ明記する。

\section{卒業までに作成する成果物}

日付ごとの細かな作業予定は示さず、卒業時に確認できる成果物を表\ref{tab:deliverables}にまとめる。

\begin{table*}[t]
  \centering
  \setlength{\tabcolsep}{4pt}
  \caption{卒業までに作成する成果物}
  \label{tab:deliverables}
  \begin{tabular}{L{0.22\textwidth}L{0.38\textwidth}L{0.29\textwidth}}
    \toprule
    成果物 & 含める内容 & 完了の確認 \\
    \midrule
    院内比較用システム & 専門家が確認したタスクグラフ、提案UI、現行電子カルテと条件をそろえた課題 & 院内サーバーで起動し、課題と必要情報を確認できる \\
    技術評価の記録 & 8種類の不整合、4種類の検査方法、正式結果、再現に必要な設定と集計 & 保存済み成果物から主要指標を再計算できる \\
    ユーザー評価の記録 & 研究計画、倫理・実施許可の確認記録、実施した場合の匿名化した操作指標、質問紙、分析、聞き取りの要約 & 実施条件と結果、または未実施の理由と残る課題を報告できる \\
    修士論文と発表資料 & 認知負荷の問題、タスクグラフ、院内実装、技術評価、ユーザー評価を結ぶ研究全体 & 主張と各評価の範囲を分けて説明できる \\
    \bottomrule
  \end{tabular}
\end{table*}

修士論文では、タスクグラフによるUI生成を提案の中心に置く。整合性評価は、生成した画面を利用者へ提示する前に要求と表示の対応を確かめる技術評価としてまとめる。ユーザー評価は、提案UIが現行電子カルテでの情報探索と統合の負担を減らすかを検証する評価として報告する。技術評価の結果だけから認知負荷の軽減を主張せず、それぞれの結果を対応する研究質問へ結び付ける。

\section{制約}

現時点で認知負荷の軽減は未評価であり、院内で動作したことだけから利用効果を結論付けられない。ユーザー評価の診療場面、参加者、症例数が限られる場合、結果を他の診療科や施設へそのまま一般化できない。提案UIは情報確認を支援するものであり、診断や治療方針を自動で決めるものではない。

タスクグラフが表す必要情報と表示方法には、診療専門家の確認が要る。整合性の技術評価は、決められた要求が表示まで保たれるかを検査するが、要求自体が医学的に正しいことは保証しない。この二つを分け、専門家による内容確認、機械的な整合性検査、利用者による比較評価を順に行う。

\section{まとめ}

本研究は、電子カルテに分散した情報を探して統合する認知負荷を軽減するため、診療タスクに合わせたUIを生成する。Scenario、Task、Data、Widgetを結ぶタスクグラフを用い、院内データベースに対して読み取り専用SQLを実行し、対話的な画面を構成する。システム面では、現行電子カルテと同じ院内環境で比較できる状態にある。

完了した層間整合性評価では、タスクから表示までの関係をまとめて検査する二方式が、763件の不整合をすべて拒否し、妥当な更新をすべて受理した。この結果はユーザー評価の技術的な前提であり、認知負荷の軽減を示す結果ではない。必要な確認と許可が整えば、卒業までに、専門家が確認した課題を使って現行電子カルテと提案UIを比較し、正確さ、時間、操作、NASA-TLX、SUS、聞き取りから利用効果を評価する。

\bibliographystyle{junsrt}
\bibliography{references}

\end{document}
